\documentclass[11pt]{article}
\usepackage{amsmath,amssymb,amsthm}
\newtheorem{theorem}{Theorem}
\newcommand{\R}{\mathbb R}
\newcommand{\ket}[1]{\lvert #1\rangle}
\usepackage[colorlinks]{hyperref}
\usepackage{oo-journal}
\ooshorttitle{Phase-invariant light}
\title{A sharp Wigner--Husimi bound for phase-invariant light}
\hypersetup{pdfauthor={GPT-6.2 Astra (high), Aryan Singh},pdftitle={A sharp Wigner-Husimi bound for phase-invariant light}}
\author{GPT-6.2 Astra (high)\enspace{\small(lead AI author of this edition)}\\\href{https://www.researchgate.net/profile/Aryan-Singh-101}{Aryan Singh}\enspace{\small$\cdot$\enspace oo research\enspace$\cdot$\enspace The Open University (student)}}
\date{26 November 2026}
\begin{document}
\maketitle
\begin{abstract}
Gagatsos conjectured an upper bound on single-mode Wigner negativity in terms of the Husimi distribution. We prove the conjectured inequality for every photon-number-diagonal state, with no cutoff on photon number. The vacuum is the unique state in this class attaining finite equality. The proof bounds a Laguerre absolute integral by an exactly computable fourth moment and compares that moment with the Husimi integral by induction.
\end{abstract}

\section{What can a heterodyne distribution tell us?}
Let $W_\rho$ be the Wigner function and $Q_\rho$ the probability density of an ideal heterodyne measurement on one optical mode. We use $[x,p]=i$, $z=(x,p)$, $\alpha=(x+ip)/\sqrt{2}$ and $Q_\rho(z)=\langle\alpha|\rho|\alpha\rangle/(2\pi)$. Both distributions have integral one in $d^2z=dx\,dp$. If $W_\rho$ is integrable, its negative volume is $(\|W_\rho\|_1-1)/2$. Define
\[
 H(\rho)=\frac{1}{8\pi}\left(\int_{\R^2}\sqrt{Q_\rho(z)}\,d^2z\right)^2.
\]
Gagatsos conjectured $\|W_\rho\|_1\le H(\rho)$ for every single-mode state and proved $\|W_\rho\|_1\le 2H(\rho)$ in general \cite{gagatsos}. States with nonnegative Wigner functions already obey the conjecture by the Lieb--Solovej bound cited there. Excited photon-number states have negative Wigner regions \cite{kenfack}.

\begin{theorem}
For every state invariant under all optical phase rotations, $\|W_\rho\|_1\le H(\rho)$ in the extended nonnegative reals. When $H(\rho)<\infty$, equality holds exactly for the vacuum. The constant $1/(8\pi)$ is sharp.
\end{theorem}
Phase rotation acts on $\ket n$ by $e^{-in\theta}$. Invariance for every $\theta$ forces every off-diagonal density-matrix element to vanish. Thus the class consists precisely of arbitrary probability mixtures of photon-number states.

\section{An all-orders estimate for a number state}
Write $r=|z|$, and let $L_n$ be the ordinary Laguerre polynomial with $L_n(0)=1$. The number-state densities are
\[
 W_n(z)=\frac{(-1)^n}{\pi}e^{-r^2} L_n(2r^2),\qquad
 Q_n(z)=\frac{e^{-r^2/2}}{2\pi n!}\left(\frac{r^2}{2}\right)^n.
\]
In polar coordinates substitute $u=2r^2$ in the Wigner integral and $v=r^2/2$ in the Husimi integral. The resulting quantities are
\begin{equation}\label{radial}
 A_n=\frac12\int_0^\infty e^{-u/2}|L_n(u)|\,du,\qquad
 B_n=H(\ket n\!\langle n|)=\frac{2^n\Gamma(1+n/2)^2}{n!}.
\end{equation}
In particular $A_0=B_0=1$. For $f_n(u)=e^{-u/2}L_n(u)$ the Laguerre orthogonality relation gives $\int f_n^2=1$. Let $M_n=\int_0^\infty u^4 f_n(u)^2\,du$. For $a>0$, Cauchy--Schwarz with weight $1+(u/a)^4$ gives
\[
 4A_n^2\le \left(1+\frac{M_n}{a^4}\right)
 \int_0^\infty\frac{du}{1+(u/a)^4}
 =\left(1+\frac{M_n}{a^4}\right)\frac{\pi a}{2\sqrt{2}} .
\]
Here the last integral follows from the beta integral. Taking $a=(3M_n)^{1/4}$ and raising to the fourth power yields
\begin{equation}\label{bound}
 A_n^8\le \frac{\pi^4}{1728}M_n .
\end{equation}

To evaluate $M_n$, apply the three-term recurrence \cite{dlmf}
\[
 uL_j=-(j+1)L_{j+1}+(2j+1)L_j-jL_{j-1}
\]
twice and take the squared norm of $u^2L_n$ in $L^2(e^{-u}du)$. In order of Laguerre indices $n-2,\ldots,n+2$, the coefficients are
\[
 n(n-1),\quad -4n^2,\quad 6n^2+6n+2,\quad -4(n+1)^2,\quad(n+1)(n+2).
\]
Negative indices are omitted and their coefficients vanish. Summing squares gives
\[
 M_n=70n^4+140n^3+170n^2+100n+24.
\]
The gamma recurrence in \eqref{radial} gives $B_{n+2}/B_n=(n+2)/(n+1)$. Consequently $B_{n+2}^8/M_{n+2}>B_n^8/M_n$ whenever
\[
 D(n)=(n+2)^8 M_n-(n+1)^8 M_{n+2}>0.
\]
The coefficients of $D$ in increasing powers of $n$ are
\[
 (3000,20324,60606,102244,107730,74088,33516,9672,1620,120).
\]
Every coefficient is positive. For the odd base, $M_1=504$, $B_1=\pi/2$ and $1728\cdot3^4>2^8\cdot504$; hence $B_1^8>\pi^4 M_1/1728$. For the even base, $M_2=3144$, $B_2=2$ and $1728\cdot 2^8>100\cdot3144$. Here $3<\pi<22/7$ implies $\pi^4<100$. Induction on each parity and \eqref{bound} give $A_n<B_n$ for every $n\ge 1$.

\section{Arbitrary photon-number statistics}
Let $\rho=\sum_{n\ge0}p_n\ket n\!\langle n|$. Both phase-space functions are linear in $\rho$; pointwise convergence follows from trace-norm convergence and the bounded displaced parity and coherent projectors defining them. The triangle inequality and Tonelli's theorem give $\|W_\rho\|_1\le\sum_n p_n A_n$. To bound the other side, apply the integral triangle inequality for the Euclidean norm to the vector with components $\sqrt{p_nQ_n(z)}$, $0\le n\le N$. This gives
\[
 \left(\sum_{n=0}^N p_n\left[\int \sqrt{Q_n}\right]^2\right)^{1/2}
 \le \int\sqrt{\sum_{n=0}^N p_nQ_n}\le\int\sqrt{Q_\rho}.
\]
Taking $N\to\infty$ and squaring proves
\[
 \|W_\rho\|_1\le\sum_n p_n A_n\le\sum_n p_n B_n\le H(\rho).
\]
If $H(\rho)$ is finite and some $p_n>0$ for $n\ge1$, the middle inequality is strict. The vacuum saturates it as well as both endpoint bounds. This proves the theorem.

\paragraph{Physical check.}
For one photon, $L_1(u)=1-u$. Splitting \eqref{radial} at $u=1$ gives $A_1=4e^{-1/2}-1$, whereas $B_1=\pi/2$. The Wigner negative volume is $2e^{-1/2}-1\approx0.2131$ and the heterodyne upper bound is $(\pi/2-1)/2\approx0.2854$. This one-photon value is known \cite{kenfack}. The theorem establishes the bound for every photon count and every probability distribution over counts. It leaves the general question for photon-number coherences open. Applying the same displacement to the whole state simply translates both integrals.

\paragraph{Edition note.} For a diagonal mixture, the probabilities appear linearly in the upper bound on the Wigner norm. The integral triangle inequality puts the same weighted sum of number-state Husimi bounds below the Husimi functional of the mixture. This step admits infinite photon-number support and uses the absence of off-diagonal terms in the number basis.

\begin{thebibliography}{9}
\bibitem{gagatsos} C. N. Gagatsos, \emph{Wigner negativity costs phase-space area}, Conjecture 1, Eq. (7), and Propositions 1--2 (2026). \href{https://arxiv.org/abs/2611.30969}{arXiv:2611.30969v1}.
\bibitem{kenfack} A. Kenfack and K. {\.{Z}}yczkowski, \emph{Negativity of the Wigner function as an indicator of nonclassicality}, J. Opt. B \textbf{6}, 396--404 (2004). \href{https://arxiv.org/abs/quant-ph/0406015}{arXiv:quant-ph/0406015}.
\bibitem{dlmf} NIST Digital Library of Mathematical Functions, \href{https://dlmf.nist.gov/18.3}{\S18.3} and \href{https://dlmf.nist.gov/18.9}{\S18.9}, Laguerre orthogonality and recurrence.
\end{thebibliography}
\end{document}
